\documentclass[10pt,a4paper]{amsart}
\usepackage[margin=19mm]{geometry}
\usepackage{amsmath,amssymb,mathtools}
\usepackage[scr=rsfs,cal=euler]{mathalfa}
\usepackage{xfrac}
\usepackage[unicode,hidelinks,bookmarks=false]{hyperref}
\newcommand{\ie}{i.e.\ }
\newcommand{\cf}{cf.\ }
\newcommand{\resp}{respectively\ }
\newcommand{\Subsection}{\subsection}
\providecommand{\nobreakdash}{\nobreak\hbox{-}}
\setlength{\emergencystretch}{2em}
\multlinegap0pt
\usepackage{mathtools}
\usepackage{mathrsfs}
\usepackage{enumitem}
\numberwithin{equation}{section}
\newtheorem{theorem}{Theorem}[section]
\newtheorem{conjecture}{Conjecture}[section]
\newtheorem{proposition}[theorem]{Proposition}
\newtheorem{lemma}[theorem]{Lemma}
\newtheorem{corollary}[theorem]{Corollary}
\newtheorem{remark}[theorem]{Remark}
\DeclareMathOperator{\Ric}{Ric}
\DeclareMathOperator{\Riem}{Riem}
\DeclareMathOperator{\tr}{tr}
\newcommand{\Hess}{\nabla^2}
\newcommand{\dd}{\,d}
\newcommand{\dv}{\,dV}
\newcommand{\ip}[2]{\left\langle #1,#2\right\rangle}
\newcommand{\norm}[1]{\left|#1\right|}
\newcommand{\Deltaf}{\Delta_f}
\begin{document}
\title[Scalar-Flatness for Critical Metrics of the L^2-Scalar Curva]{Scalar-Flatness for Critical Metrics of the $L^2$-Scalar Curvature Functional in Dimensions $5\le n\le 9$}
\author{Heng Zhang}
\date{}
\maketitle
\noindent\emph{Source and licence.} Scalar-Flatness for Critical Metrics of the $L^2$-Scalar Curvature Functional in Dimensions $5\le n\le 9$. Version arXiv:2606.28897v1 (CC BY 4.0). This material is distributed under the Creative Commons Attribution 4.0 International licence, http://creativecommons.org/licenses/by/4.0/, as recorded in the arXiv OAI licence metadata for this exact version and shown on the abstract page https://arxiv.org/abs/2606.28897v1. Full rights notice: licenses/arxiv-2606.28897-CC-BY-4.0.txt.\par\smallskip
\noindent\textit{Editorial scope.} Counted unit: the original \texttt{lemma} environment \texttt{lem:log} at source lines 393--414 together with its complete proof at lines 416--514; the delivered document keeps source lines 389--515 so that every referenced label and every citation resolves inside it. Native editable source is the paper's own concatenated TeX; the AMS wrapper is editorial formatting only. No publisher class, upstream script or unrelated artwork is redistributed.
\section{A Logarithmic growth lemma}\label{sec:log}

The next elementary lemma is useful in the completeness argument.


\begin{lemma}[Logarithmic growth lemma]\label{lem:log}
Let $v>0$ be a $C^2$ function on $(0,T)$, where $0<T<\infty$.  Assume that there exist constants $\rho>0$ and $b\geq 0$ such that
\begin{equation}\label{log-assumption}
        \rho\int_0^T \frac{v''}{v}\phi^2\,dt
        \leq
        \int_0^T (\phi')^2\,dt+b\int_0^T \phi^2\,dt
\end{equation}
for every $\phi\in C_c^\infty(0,T)$.  Set
\begin{equation}\label{prho}
        p_\rho=\frac{\sqrt{1+\rho^{-1}}-1}{2},
\end{equation}
so that $\rho p_\rho(p_\rho+1)=1/4$.  Then
\begin{equation}\label{growth-conclusion}
        \limsup_{t\to T^-}
        \frac{\log v(t)}{\log \frac1{T-t}}
        \leq p_\rho .
\end{equation}
Consequently, for every $a>0$ such that $a p_\rho<1$, one has
\begin{equation}\label{integrability-conclusion}
        \int^{T} v(t)^a\,dt<\infty .
\end{equation}
\end{lemma}

\begin{proof}
Put $x=T-t$ and $V(x)=v(T-x)$.  Let
$$
        y=(\log V)' .
$$
Then $v''(t)/v(t)=V''(x)/V(x)=y'+y^2$.  Rewriting \eqref{log-assumption} in the $x$ variable gives
$$
        \rho\int_0^T (y'+y^2)\phi^2\,dx
        \leq
        \int_0^T (\phi')^2\,dx+b\int_0^T\phi^2\,dx .
$$
For compactly supported smooth $\phi$,
\begin{equation}\label{square-identity}
        \int (y'+y^2)\phi^2
        =\int (y\phi-\phi')^2-\int (\phi')^2 .
\end{equation}
Hence
\begin{equation}\label{square-bound}
        \rho\int_0^T (y\phi-\phi')^2\,dx
        \leq
        (1+\rho)\int_0^T(\phi')^2\,dx+b\int_0^T\phi^2\,dx .
\end{equation}

Fix $R\in(0,T)$ sufficiently small.  Given $0<r<R/8$, choose a smooth cutoff $\chi_{r,R}$ such that
$$
\begin{array}{ll}
        \chi_{r,R}=1 &\text{on }[2r,R/2],\\
        \chi_{r,R}=0 &\text{near }0\text{ and near }R,
\end{array}
$$
with
$$
        |\chi_{r,R}'|\leq C/r\quad\text{on }[r,2r],
        \qquad
        |\chi_{r,R}'|\leq C/R\quad\text{on }[R/2,R],
$$
and analogous second-derivative bounds.  The function
$$
        \phi(x)=x^{1/2}\chi_{r,R}(x)
$$
is smooth and compactly supported in $(0,T)$, because $\chi_{r,R}$ is identically zero near $x=0$.  Equivalently, one may start from the ideal piecewise-smooth cutoff and approximate it in $H_0^1(0,T)$.  Substitution in \eqref{square-bound} yields
$$
        \int_0^T(\phi')^2\,dx
        =\frac14\log\frac{R}{r}+O(1),
        \qquad
        \int_0^T\phi^2\,dx=O(R^2),
$$
where $R$ is fixed and the constants in $O(1)$ are independent of $r$.  Since $R$ is fixed, the term $bO(R^2)$ is also independent of $r$ and disappears after division by $\log(1/r)$ and passage to the limit $r\to0^+$.  On $[2r,R/2]$,
$$
        y\phi-\phi'=x^{1/2}\left(y-\frac1{2x}\right).
$$
Therefore
\begin{equation}\label{central-bound}
        \int_{2r}^{R/2} x\left(y-\frac1{2x}\right)^2\,dx
        \leq
        \frac{1+\rho}{4\rho}\log\frac{R}{r}+O(1).
\end{equation}
By Cauchy--Schwarz,
\begin{align*}
&\log V(2r)-\log V(R/2)+\frac12\log\frac{R/2}{2r} \\
&\qquad
= -\int_{2r}^{R/2}\left(y-\frac1{2x}\right)\,dx \\
&\qquad
\leq
\left|\int_{2r}^{R/2}\left(y-\frac1{2x}\right)\,dx\right| \\
&\qquad
\leq
\left(\int_{2r}^{R/2}x\left(y-\frac1{2x}\right)^2\,dx\right)^{1/2}
\left(\int_{2r}^{R/2}\frac{dx}{x}\right)^{1/2} \\
&\qquad
\leq
\sqrt{\frac{1+\rho}{4\rho}}\log\frac{R/2}{2r}+o\!\left(\log\frac1r\right)
\end{align*}
as $r\to0^+$.  Dividing by $\log(1/r)$ and letting $r\to0^+$ first gives
$$
        \limsup_{r\to0^+}\frac{\log V(2r)}{\log(1/r)}
        \leq
        \sqrt{\frac{1+\rho}{4\rho}}-\frac12
        =p_\rho .
$$
Replacing $2r$ by $s$ does not change the limsup, since
$\log(1/r)=\log(1/s)+\log 2$ and hence
$\log(1/r)/\log(1/s)\to1$ as $s\to0^+$.  Therefore
$$
        \limsup_{s\to0^+}\frac{\log V(s)}{\log(1/s)}
        \leq p_\rho .
$$
This is \eqref{growth-conclusion}.  If $a>0$ and $a p_\rho<1$, choose $\varepsilon>0$ with $a(p_\rho+\varepsilon)<1$.  From \eqref{growth-conclusion}, after decreasing the terminal interval if necessary, one has
$$
        v(t)\leq C_\varepsilon (T-t)^{-p_\rho-\varepsilon}
        \qquad\text{for }t\text{ sufficiently close to }T.
$$
Thus
$$
        \int^T v(t)^a\,dt
        \leq
        C\int^T (T-t)^{-a(p_\rho+\varepsilon)}\,dt<\infty .
$$
\end{proof}


\end{document}
